\documentclass[11pt,a4paper]{article}
\usepackage{iftex}
\ifPDFTeX
  \usepackage[T1]{fontenc}
  \usepackage[utf8]{inputenc}
  \usepackage{lmodern}
  \input{glyphtounicode}
  \pdfgentounicode=1
\else
  \usepackage{fontspec}
  % Kerning off: kerned pairs ("A WS", "T ree") split words for ATS text extraction.
  \setmainfont{lmroman10-regular}[Extension=.otf, BoldFont=lmroman10-bold,
    ItalicFont=lmroman10-italic, BoldItalicFont=lmroman10-bolditalic, Kerning=Off]
\fi
\usepackage[margin=0.9in]{geometry}
\usepackage[hidelinks]{hyperref}
\hypersetup{pdftitle={Abhinav Kaushik - Cover Letter - openai}, pdfauthor={Abhinav Kaushik}}

\pagestyle{empty}
\setlength{\parindent}{0pt}
\setlength{\parskip}{10pt}
\raggedright

\begin{document}

{\Large\bfseries Abhinav Kaushik}\\
Applied AI Engineer\\
New Delhi, India \textbar{} +91 8700571859 \textbar{} \href{mailto:aabhinav.kaushik@gmail.com}{aabhinav.kaushik@gmail.com} \textbar{} \href{https://www.linkedin.com/in/abhinav-kaushik10/}{LinkedIn}  

September 27, 2026

Hiring Team\\
openai

\textbf{Re: Applied AI Engineer, Startups (Codex)}

Dear Hiring Team,

I am applying for the Applied AI Engineer role on the Startups team in Munich, working with Codex. Much of my career has been spent building AI applications side by side with the people who need them, then making them dependable enough to stay in daily use. Getting a good idea into a workflow that sticks is the common thread.

Codex is reshaping how young companies build software, and Munich's startup scene is a place where that change is happening quickly. I want to help founders there get lasting value from it. My fiancée lives in Berlin, and we plan to move to Munich together and settle there long term.

I would bring a builder's habits to the role. At Viavi I prototyped 15+ AI applications from client requirements, and 7 of them moved into production or paid client work. At United Health Group I shipped a lab-extraction agent that automated the work behind 800 Jira tickets, with evaluation pipelines to keep it reliable.

Outside work I am building ThinkTree.AI, a non-profit AI learning platform used by NGO teachers to support students with ADHD. It keeps me asking whether a tool truly helps the people using it, a question I would bring to every startup I work with. I would be glad to talk further.

Sincerely,\\[4pt]
Abhinav Kaushik

\end{document}